\section{Research Questions}
\label{sec:research-questions}

The survey is guided by six research questions designed to move from descriptive mapping to integration analysis and gap identification.

\begin{description}

\item[RQ1:] \textbf{Which quality dimensions of LLM-based systems are evaluated in the existing literature?}

RQ1 maps the literature onto the lifecycle-oriented taxonomy and examines how attention is distributed across input/specification, context/retrieval, output, conversation/interaction, agent/process/trajectory, and operational/system quality. It also records cross-cutting concerns such as robustness, safety, security, fairness, multimodality, and efficiency.

\item[RQ2:] \textbf{Which evaluation methods and protocols are used to assess these quality dimensions?}

RQ2 separates the construct being evaluated from the mechanism used to measure it. We examine deterministic and rule-based checks, lexical and reference-based metrics, semantic metrics, learned evaluators, human evaluation, LLM-as-a-Judge, multi-judge approaches, task- and environment-based verification, adversarial testing, perturbation-based methods, trace-aware evaluation, and statistical reliability analysis. We additionally record protocols such as absolute scoring, pairwise comparison, ranking, rubric scoring, binary classification, critique, preference judgment, and task-success verification.

\item[RQ3:] \textbf{To what extent does the literature connect quality across multiple stages of an LLM-based system?}

RQ3 investigates whether evaluation remains isolated at one component, explicitly relates multiple stages, measures a multi-stage workflow only through its final outcome, or provides a lifecycle-integrated assessment. This question directly addresses the difference between black-box end-to-end evaluation and stage-aware quality attribution.

\item[RQ4:] \textbf{How is the reliability of the evaluation mechanism itself assessed?}

RQ4 treats evaluator quality as a meta-evaluation problem. We examine whether studies assess consistency, human agreement, calibration, bias, robustness, prompt sensitivity, dependence among multiple evaluators, or other properties that determine whether an automated or human judgment can be trusted.

\item[RQ5:] \textbf{How are operational and deployment-oriented properties incorporated into LLM quality assessment?}

RQ5 examines whether semantic or task-level quality is considered together with latency, throughput, monetary cost, token and memory consumption, energy use, availability, observability, privacy, safety, and related system-level constraints. The goal is to identify whether operational quality is treated as part of overall system quality or as a separate engineering concern.

\item[RQ6:] \textbf{What evidence gaps remain toward reliable end-to-end quality assessment of LLM-based systems?}

RQ6 synthesizes the preceding results to identify dimensions that are well studied in isolation but weakly connected, evaluation methods whose reliability is insufficiently characterized, and lifecycle relationships for which empirical evidence remains scarce. This final question provides the basis for the research agenda proposed in the survey.

\end{description}
